\documentclass[11pt,a4paper]{article}

\usepackage[margin=2.5cm]{geometry}
\usepackage[T1]{fontenc}
\usepackage{lmodern}
\usepackage{amsmath,amssymb}
\usepackage{booktabs}
\usepackage{graphicx}
\usepackage{subcaption}
\usepackage{xcolor}
\usepackage[colorlinks=true,linkcolor=blue!50!black,urlcolor=blue!50!black]{hyperref}

% Highlighted placeholder for things still to write or verify
\newcommand{\todo}[1]{\textcolor{red}{[TODO: #1]}}

\title{Image Dithering: Notes}
\author{\todo{Your name}}
\date{\today}

\begin{document}
\maketitle
\tableofcontents
\bigskip

% ------------------------------------------------------------------
\section{What is Dithering?}
\label{sec:what}

A digital image stores each pixel as a value from a limited set of
intensities. When the display (or file format) can only represent a
\emph{reduced} palette, every pixel must be mapped to one of the available
values. \emph{Dithering} is the family of techniques that arranges the
available values spatially so that, viewed from a distance, the eye
averages neighbouring pixels and perceives intermediate tones that the
palette cannot represent directly.

\subsection{Quantisation}

Let a single channel of an image be a function $I(x,y)\in[0,255]$.
A quantiser with $L$ levels snaps each value to the nearest of $L$ evenly
spaced values:
\begin{equation}
  \Delta = \frac{255}{L-1}, \qquad
  q(v) = \Delta \cdot \operatorname{round}\!\left(\frac{v}{\Delta}\right).
\end{equation}
The \emph{quantisation error} of a pixel is $e = v - q(v)$, and always
satisfies $|e|\le \Delta/2$. For a colour image the quantiser is applied to
the red, green and blue channels independently, giving $L^3$ possible colours
(Table~\ref{tab:levels}).

\begin{table}[htbp]
  \centering
  \caption{Palette size as a function of levels per channel $L$.}
  \label{tab:levels}
  \begin{tabular}{rrrr}
    \toprule
    $L$ & Grayscale shades & Colours ($L^3$) & Step $\Delta$ \\
    \midrule
    2   & 2   & 8          & 255   \\
    3   & 3   & 27         & 127.5 \\
    4   & 4   & 64         & 85    \\
    6   & 6   & 216        & 51    \\
    256 & 256 & 16,777,216 & 1     \\
    \bottomrule
  \end{tabular}
\end{table}

% ------------------------------------------------------------------
\section{Historical Context}
\label{sec:history}

Early display hardware could show only a handful of colours at once. The
IBM Colour Graphics Adapter (CGA), introduced in 1981 with the IBM PC, is the
classic example. Its 16-colour palette is built from four on/off signals:
red, green, blue and an intensity bit (RGBI). Sixteen colours need only
four bits per pixel, since $2^4 = 16$.

This should not be confused with ``16-bit colour'' (High Color), which uses
16 bits per pixel and yields $2^{16} = 65{,}536$ colours
(Table~\ref{tab:bits}).

\begin{table}[htbp]
  \centering
  \caption{Bits per pixel versus number of distinct colours.}
  \label{tab:bits}
  \begin{tabular}{lrr}
    \toprule
    Name & Bits per pixel & Colours \\
    \midrule
    CGA palette          & 4  & 16         \\
    ``16-bit'' High Color & 16 & 65,536     \\
    24-bit True Color    & 24 & 16,777,216 \\
    \bottomrule
  \end{tabular}
\end{table}

In CGA's graphics modes, only four of the sixteen colours could be shown
simultaneously, chosen from fixed palettes. Programmers therefore mixed
pixel patterns to imitate the missing colours, which is a major reason
dithering became widespread. \todo{Add references for CGA specifications.}

% ------------------------------------------------------------------
\section{Use Cases: Early Game Graphics}
\label{sec:usecases}

Console and arcade artists faced the same constraint: a small palette and low
resolution. The Sega Mega Drive (Genesis) offered a 9-bit palette of 512
colours, of which only a limited number could be shown at once. Artists in
games such as \emph{Sonic the Hedgehog} used checkerboard patterns of two
colours; on a CRT television connected by composite video, the signal
blurred the pattern into an apparent third colour or a semi-transparent
effect (a commonly cited example is the waterfalls of Green Hill Zone).
\todo{Verify details and cite a source.}

Dithering also served monochrome displays: Atkinson's algorithm
(Section~\ref{sec:atkinson}) was written for the original Apple Macintosh,
whose screen was strictly one bit per pixel.

% ------------------------------------------------------------------
\section{The Primitive Approach: Nearest Colour}
\label{sec:nearest}

The simplest strategy maps every pixel independently to the closest palette
colour. With palette $P$ and pixel $p=(r,g,b)$,
\begin{equation}
  \operatorname{map}(p) = \arg\min_{c\in P}
  \sqrt{(r-c_r)^2 + (g-c_g)^2 + (b-c_b)^2}.
\end{equation}
For two grayscale levels this reduces to thresholding at 128.

\paragraph{Remark (choice of distance).}
Euclidean distance is used rather than cosine similarity. Cosine similarity
compares only the \emph{direction} of two colour vectors and ignores their
magnitude, so $(255,0,0)$ and $(128,0,0)$ would be considered identical even
though one is much darker.

Nearest-colour mapping discards the error entirely. Smooth gradients collapse
into flat regions with hard edges (\emph{banding} or \emph{contouring}), and
large areas of similar tone become a single solid colour. The algorithms in
Section~\ref{sec:algorithms} fix this by \emph{remembering} the error.
\todo{Insert a threshold-versus-dithered comparison figure.}

% ------------------------------------------------------------------
\section{Dithering Algorithms}
\label{sec:algorithms}

\subsection{Error diffusion}

All four algorithms below are \emph{error-diffusion} methods. Pixels are
processed in scan order. For the current pixel at $(x,y)$:
\begin{align}
  \text{old} &= I(x,y), \qquad \text{new} = q(\text{old}), \qquad
  e = \text{old} - \text{new}, \\
  I(x+d_x,\,y+d_y) &\mathrel{+}= \frac{w}{D}\, e
  \quad \text{for each } (d_x,d_y,w) \text{ in the kernel.}
\end{align}
The kernel lists which not-yet-processed neighbours receive a share of the
error, with integer weights $w$ and a divisor $D$. In the matrices below,
$\ast$ marks the current pixel and $\cdot$ marks pixels that are never
modified.

Our implementation uses \emph{serpentine} scanning: the scan direction
alternates every row (and the kernel is mirrored accordingly), which avoids
the diagonal ``worm'' artefacts of always scanning left to right.

\subsection{Floyd--Steinberg}
\label{sec:fs}

Published by Floyd and Steinberg in 1976, this is the most widely used
algorithm. It spreads the full error over four neighbours:
\begin{equation}
  \frac{1}{16}
  \begin{pmatrix}
    \cdot & \ast & 7 \\
    3 & 5 & 1
  \end{pmatrix}
\end{equation}

\subsection{Atkinson}
\label{sec:atkinson}

Developed by Bill Atkinson at Apple for the original Macintosh, this kernel
spreads the error over six neighbours but divides by 8, so only
$6/8 = 75\%$ of the error is propagated:
\begin{equation}
  \frac{1}{8}
  \begin{pmatrix}
    \cdot & \ast & 1 & 1 \\
    1 & 1 & 1 & \cdot \\
    \cdot & 1 & \cdot & \cdot
  \end{pmatrix}
\end{equation}
Discarding a quarter of the error gives higher contrast: highlights and
shadows tend to saturate to pure white and black, while mid-tones keep
crisp, less noisy detail.

\subsection{Jarvis--Judice--Ninke}
\label{sec:jarvis}

This kernel spreads the error over twelve neighbours across three rows:
\begin{equation}
  \frac{1}{48}
  \begin{pmatrix}
    \cdot & \cdot & \ast & 7 & 5 \\
    3 & 5 & 7 & 5 & 3 \\
    1 & 3 & 5 & 3 & 1
  \end{pmatrix}
\end{equation}
The wider spread produces smoother results with finer texture, at the cost of
more computation per pixel.

\subsection{Sierra}
\label{sec:sierra}

The three-row Sierra kernel is a lighter relative of Jarvis--Judice--Ninke,
using ten neighbours and a power-of-two divisor:
\begin{equation}
  \frac{1}{32}
  \begin{pmatrix}
    \cdot & \cdot & \ast & 5 & 3 \\
    2 & 4 & 5 & 4 & 2 \\
    \cdot & 2 & 3 & 2 & \cdot
  \end{pmatrix}
\end{equation}

\subsection{Comparison}

\begin{table}[htbp]
  \centering
  \caption{Summary of the error-diffusion kernels used.}
  \label{tab:kernels}
  \begin{tabular}{lrrr}
    \toprule
    Algorithm & Neighbours & Divisor $D$ & Error propagated \\
    \midrule
    Floyd--Steinberg        & 4  & 16 & 100\% \\
    Atkinson                & 6  & 8  & 75\%  \\
    Jarvis--Judice--Ninke   & 12 & 48 & 100\% \\
    Sierra                  & 10 & 32 & 100\% \\
    \bottomrule
  \end{tabular}
\end{table}

\todo{Verify publication years and add references for each algorithm.}
\todo{Add side-by-side result figures for the same source image.}

% ------------------------------------------------------------------
\section{Conclusion}
\label{sec:conclusion}

\todo{Summarise your findings. Observations from the analysis tool that may
be worth confirming on real images: (1) a dithered image's histogram collapses
onto $L$ values per channel, yet it still looks smooth because the tone is
carried by the \emph{density} of the dots; (2) error diffusion preserves the
mean brightness almost exactly; (3) pixelising before dithering makes the
pattern visible and keeps computation fast.}

% ------------------------------------------------------------------
\appendix
\section{Dithered Image Gallery}
\label{app:gallery}

\IfFileExists{appendix_gallery.tex}{\input{appendix_gallery.tex}}{%
  \todo{Run \texttt{generate\_gallery.py} to create \texttt{appendix\_gallery.tex}.}}

\end{document}